% !TeX root = ../main.tex

\chapter{State of the Art}\label{chapter:state-of-art}

This chapter presents the research topic of this thesis, including the building energy management in distributed networks, the existing control methods, and the remaining research gaps. The focus of this thesis is the coordinated energy management of multiple prosumers in a low-voltage distribution network. Each prosumer consists of a household load, PV generation, a BESS, and residential EV charging demand. The controller should balance the dynamic electricity price, renewable energy fluctuations, EV departure SoC requirements, and grid safety constraints. Therefore, this chapter first introduces the fundamental of HEMS, then reviews the existing control methods, and finally summarizes the limitations and challenges of these methods, thereby motivating the research presented in this thesis.

\section{Buidling Energy Management in Distributed Networks}

With the development of smart grids, advanced infrastructure and more residential distributed energy resources, household users have gradually evolved from traditional passive loads into prosumers that can more actively participate in distribution network operation. Smart HEMS is a key implementation of residential demand-side management. It can monitor and coordinate household loads, PV systems, BESS, EVs, and other energy resources in the real time, thereby reducing electricity costs and improving energy efficiency. Nowadays,HEMS has gradually evolved from simple load scheduling under the traditional centralized power supply framework toward cyber-physical HEMS with communication, sensing, control, and user response capabilities~\parencite{ZHOU201630}. Besides, the widespread deployment of residential air conditioners and other devices has made HEMS become a dynamic multi-objective optimization problem, so that it can simultaneously balance electricity cost, thermal comfort and other objectives. Since users' demand of "optimal" energy usage is very subjective and time-varying, HEMS need to adjust its control objectives according to user preferences~\parencite{DU2026127976}.

In building energy management, PV systems, BESS, and EVs are among the most important flexible energy resources, therefore we mainly focus on these three devices in this thesis. PV generation can reduce the electricity demand purchased from the distribution network. When residential PV systems are jointly managed with household loads and other energy resources, their economic value can be improved under market mechanisms. But they are inevitably affected by the intermittency of PV generation and the uncertainty of load demand~\parencite{CLASTRES201055}. Besides due to its fluctuation, it can also cause reverse power flow and voltage problems during periods of low demand and high generation. As for batteries, Dongol et al.\ investigated MPC-based peak shaving using residential batteries in a rural low-voltage distribution network, showing that residential batteries can be used to smooth PV feed-in and load peaks, although their effectiveness depends on the availability of prediction information and the design of the control objectives~\parencite{DONGOL20181}. The integration of BESS also enables other applications. Manjunatha et al.\ demonstrated the application potential of battery energy storage systems in distributed distribution networks for peak shaving, cost-optimal control, and operation under time-varying electricity prices~\parencite{6652143}. Id Omar et al.\ further discussed stacked control of behind-the-meter BESS for both local services and distribution network services, including PV self-consumption, peak-load reduction, and balancing services, indicating that BESS has evolved beyond a household cost optimization tool and has become a flexible resource for providing grid services~\parencite{IDOMAR2026100249}.

At present, xxcept for BESS and PV, EV is also more and more important. Residential charging stations have enabled EVs to become part of HEMS, while the large-scale integration of EVs has further increased the complexity of energy management. Similar to BESSs and PV systems, EV can also be regarded as flexible resources, and their charging schedules can be adjusted according to electricity prices and user mobility requirements. Unlike BESSs and PV systems, however, EVs are characterized by the connection periods, charging power limits, and departure SoC requirements. Fast EV charging stations may also have an impact on residential distribution networks. For example, high-power EV charging can lead to voltage fluctuations, harmonics, power quality degradation, and distribution network stability issues~\parencite{9336258}. Seal et al.\ investigated the integration of EVs as mobile energy storage units into a centralized MPC-based HEMS and analyzed the influence of uncertainties in EV arrival, departure, and trip discharge on the control performance through Monte Carlo simulations~\parencite{10012434}. Park and Moon proposed a MADRL-based charging scheduling method for multiple EVs in a smart grid and emphasized that decentralized execution is more suitable for practical charging environments requiring fast decision-making and heterogeneous vehicle states~\parencite{PARK2022120111}. These studies indicate that EV modeling should not simply treat EVs as conventional electrical loads, but should explicitly consider their arrival time, departure time, SoC evolution, and user requirements.

Overall, building energy management in distributed networks must simultaneously consider the economic objectives on the user side and the operational requirements on the distribution network side. Makhadmeh et al.\ pointed out that the primary objectives of the smart home power scheduling problem under dynamic electricity pricing typically include reducing electricity costs, reducing the peak-to-average ratio, and improving user comfort~\parencite{MAKHADMEH2019109362}. Ali et al.\ summarized the architecture of HEMS in smart grids from the perspective of closed-loop control, emphasizing that HEMS involves not only the scheduling of household devices but also the coordinated operation with demand-side management, demand response, communication networks, and monitoring devices~\parencite{ALI2022104609}. Therefore, in the low-voltage distribution network considered in this thesis, energy management must not only schedule the energy resources within individual households but also coordinate the simultaneous actions of multiple prosumers and their combined impacts on the net load.

\section{Existing Control Methods}
The objective functions, device models, user comfort requirements, prediction uncertainty, and computational complexity considered in HEMS research vary considerably~\parencite{BEAUDIN2015318}. 

Early HEMS control methods mainly use rule-based or heuristic strategies. These methods are simple to implement, only require low computational effort, and are also suitable for practical deployment. However, they lack systematic foresight regarding future electricity prices, PV generation, and load variations. Shakeri et al.\ proposed a HEMS architecture and control algorithm that schedules household devices according to the electricity price, battery status, and thermal appliance states. Their simulation results demonstrated a reduction in daily electricity costs, although the control logic still relied on manually designed rules~\parencite{SHAKERI2017154}. Costanzo et al.\ formulated smart home peak-load shaving as an online scheduling problem and showed that online scheduling can reduce peak demand. However, their approach mainly focused on appliance scheduling and did not address the more complex coordination of multiple energy storage systems and multiple prosumers~\parencite{5984354}. Wu et al.\ proposed a hierarchical online EMS for a residential microgrid with hybrid hydrogen--electricity storage, where a fuzzy logic controller was employed to perform energy management at different time scales. Although this method avoids dependence on prediction models, the authors also pointed out that the lack of foresight may prevent stable electricity purchasing during low-price periods~\parencite{WU2024123020}. Rahim et al.\ compared heuristic algorithms such as GA, BPSO, and ACO for home energy management controllers, demonstrating that heuristic methods can reduce electricity costs and the peak-to-average ratio, although their performance strongly depends on the algorithm design and problem formulation~\parencite{RAHIM2016452}. Sisodiya et al.\ applied PSO to demand response scheduling for household devices including EVs and UPS systems, demonstrating the application of metaheuristic methods in residential energy storage and load management~\parencite{8385345}.

Different from rule-based methods, optimization-based methods can provide a more explicit representation of device constraints and economic objectives. Zhang et al.\ proposed a demand response and HEM system that combines machine learning, optimization, and data structure design. Household loads are classified into three categories, namely fixed, regulatable, and deferrable loads. A learning-based energy consumption model is established for HVAC systems, and the demand response strategy is subsequently generated through optimization~\parencite{7450180}. Jin et al.\ proposed the Foresee HEMS based on a multi-objective MPC framework, simultaneously considering energy cost, thermal comfort, user convenience, and carbon emissions. Machine learning algorithms are employed to learn appliance models and usage patterns, thereby providing demand response services while satisfying user requirements~\parencite{JIN20171583}. Chen et al.\ proposed MPC-based appliance scheduling, where thermal and non-thermal appliances are optimized over a finite prediction horizon under time-varying retail electricity pricing, and new prediction information is continuously incorporated through a rolling optimization framework~\parencite{Chen20131401}. Beaudin et al.\ proposed a moving window algorithm that periodically updates forecasts to improve the robustness of residential energy management against forecast errors and external disturbances~\parencite{6465896}. Zheng et al.\ investigated load shifting and peak reduction using building energy storage from the perspective of demand response in the residential sector, demonstrating that energy storage systems can provide demand response capability without changing the actual electricity consumption behavior~\parencite{ZHENG2014STORAGE}.

The advantages of MPC are that it can explicitly consider SoC dynamics, power limits, comfort constraints, and forecast information during each rolling optimization horizon. So it is suitable for the scheduling of BESS, EVs, and PV systems. However, MPC and mixed-integer optimization also face problems including high computational complexity, strong model dependence, and reliance on optimization solvers. Mueller et al.\ compared rule-based control, MPC, imitation learning, and reinforcement learning for controlling multiple household BESS in a German low-voltage benchmark distribution network. This showed that model-based methods generally achieve the best objective values and the fewest constraint violations, although at the cost of significantly higher computational time. In contrast, model-free methods require less computational effort but may result in higher objective values and more distribution network constraint violations~\parencite{MUELLER2025126465}. Therefore, MPC is adopted as the baseline method in this thesis.

Distributed energy system management consists of a large number of small-scale distributed energy resources. Their distributed nature, intermittent behavior, and bidirectional power flow make system operation increasingly difficult to manage using traditional centralized methods. Consequently, besides the methods mentioned above, artificial intelligence and machine learning have been also widely applied to forecasting, optimization, real-time control, intelligent demand response, and decentralized coordination to improve the adaptability, scalability, and robustness of distributed energy systems~\parencite{SATPATHY2026127109}. Liu et al.\ applied DQN and DDQN to HEMS appliance scheduling and compared their performance with PSO, demonstrating that model-free DRL can reduce electricity costs through action--reward interactions in dynamic environments, although its performance is strongly influenced by reward design and generalization capability~\parencite{Liu2020572}. Mocanu et al.\ were among the first to apply DRL to online optimization for building energy management. Using the Pecan Street dataset containing PV systems, EVs, and building appliances, they evaluated Deep Q-learning and deep policy gradient methods, showing that deep policy gradient is more suitable for online energy resource scheduling~\parencite{MOCANU20193698}. Amer et al.\ systematically discussed the state--action pairs, reward functions, and hyperparameters of DRL for demand response in HEMS, pointing out that the performance of DRL is highly dependent on the configuration, while reward formulation and environmental variations significantly affect the convergence speed, robustness, and scheduling performance~\parencite{en15218235}. Shuvo and Yilmaz proposed a HERS framework that incorporates resident feedback and activity information into deep reinforcement learning, making HEMS more user-centric. However, they also noted that their model does not include PV systems, energy storage, microgrids, or a multi-agent framework, leaving room for further extension~\parencite{9733029}. Furthermore, Du et al.\ proposed an LLM-augmented reinforcement learning framework, in which an LLM is used during offline iterations to translate natural language user preferences into reward and state modifications, while continuously refining the policy through result analysis and policy warm-start to reduce the mismatch between the learned policy and user preferences~\parencite{DU2026127976}.

Moreover, in multi-user and distributed scenarios, MADRL is more suitable than single-agent DRL for modeling the coordinated decision-making of multiple households or multiple EVs. Xu et al.\ proposed a multi-agent reinforcement learning-based data-driven HEM method, formulating hour-ahead scheduling as a finite Markov decision process. Electricity prices and PV generation are predicted using ELM, followed by multi-agent Q-learning for scheduling household appliances and EV charging~\parencite{8981876}. Lee et al.\ adopted a centralized training and decentralized execution multi-agent deep actor--critic framework, allowing each household agent to rely only on local observations during execution, thereby improving both scalability and privacy preservation~\parencite{9302935}. Wang et al.\ proposed a scalable MADRL framework based on multi-stage PPO and imitation learning for residential hybrid energy systems, employing centralized training and decentralized execution to improve the scalability and robustness of multi-agent systems~\parencite{WANG2024123414}. The EV charging scheduling study by Park and Moon also adopted the CTDE framework, where a decentralized charging policy is learned for each EV in a smart grid charging station~\parencite{PARK2022120111}. These studies demonstrate that MADRL is capable of handling the coupling among multiple agents and the challenges associated with local observations, which is also the reason why this thesis adopts a multi-agent actor--critic framework.

At the same time, safety constraints and physical feasibility are essential considerations for the practical deployment of energy management systems. Li et al.\ reviewed the applications of DRL in smart grid operations and pointed out that DRL is capable of handling uncertainty, the curse of dimensionality, and the absence of accurate system models. However, challenges remain in reward design, sample efficiency, and simulation-to-reality transfer~\parencite{10241311}. Zhang et al.\ proposed an interior-point policy optimization-based MADRL framework that formulates secure HEM as a constrained Markov decision process and integrates non-stationary Transformer forecasting to improve cost efficiency, comfort, and safety guarantees under uncertain environments~\parencite{ZHANG2024124155}. Quakernack et al.\ applied DDQN-based multi-agent control to low-voltage distribution network stability by jointly controlling EV charging and battery storage system charging/discharging to mitigate voltage drops, transformer overload, and reverse power flow~\parencite{9960416}. Based on these, this thesis will also consider the safety constraints of the distribution network.


\section{Research Gaps}

In summary, the existing literature has investigated HEMS, residential microgrids, BESS scheduling, EV charging, MPC, machine learning, DRL, and MADRL from multiple perspectives. However, several research gaps still remain.

Many HEMS studies primarily focus on the internal optimization of a single household or a single building, with objectives such as reducing electricity costs, improving user comfort, or increasing PV self-consumption~\parencite{ZHOU201630,BEAUDIN2015318,JIN20171583}, but they pay limited attention to the network coupling effects arising when multiple prosumers are connected to the same low-voltage distribution network. In this thesis, the battery charging/discharging, PV curtailment, and EV charging decisions of multiple prosumers jointly modify the feeder net load and subsequently affect the bus voltage, line loading, and transformer loading through power flow calculations. Therefore, the conclusions obtained from single-household HEMS cannot be directly extended to multi-prosumer distribution network energy management.

Traditional optimization and MPC methods can explicitly represent device constraints and future forecast information, but their performance strongly depends on accurate system models, forecast quality, and optimization efficiency. MPC has demonstrated strong performance in residential appliance scheduling, battery peak shaving, and EV-as-mobile-storage problems~\parencite{Chen20131401,DONGOL20181,10012434}. The comparative study by Mueller et al.\ also showed that model-based methods generally achieve better objective values and fewer constraint violations~\parencite{MUELLER2025126465}. However, when multiple prosumers, BESSs, PV systems, EVs, dynamic electricity prices, and distribution network constraints are simultaneously considered, both the optimization problem size and computational complexity increase significantly. In particular, MISOCP and mixed-integer formulations become more difficult to deploy in real time due to binary variables and nonlinear power flow approximations.

Furthermore, although DRL and MADRL are suitable for handling uncertainty and online decision-making, training stability, reward design, and constraint satisfaction remain key challenges. Existing DRL-based HEMS studies have shown that DQN, DDQN, deep policy gradient, and multi-agent Q-learning can reduce residential energy costs or improve demand response performance~\parencite{Liu2020572,MOCANU20193698,8981876}. However, Amer et al.\ pointed out that the performance of DRL is highly dependent on the reward function and hyperparameters~\parencite{en15218235}. Li et al.\ also emphasized that reward design, exploration, and safety remain important limitations in DRL applications for smart grids~\parencite{10241311}. Therefore, in this thesis, different reward designs and constraint handling methods are investigated to analyze their effects on economic performance, EV feasibility, and distribution network safety. 

Although existing multi-agent methods are used to support decentralized execution for multiple users or multiple EVs, studies that simultaneously consider residential BESS, PV utilization, EV departure SoC, forecasted electricity price, load and PV generation, together with distribution network operational indicators, remain limited. Xu et al.\ and Lee et al.\ demonstrated the potential of MARL for HEM and demand-side scheduling~\parencite{8981876,9302935}. Park and Moon showed that the CTDE framework is well suited for multi-EV charging scheduling~\parencite{PARK2022120111}, while Wang et al.\ demonstrated that MADRL can improve the scalability and training efficiency of residential hybrid energy systems~\parencite{WANG2024123414}. However, these studies mainly focus on charging stations, building thermal comfort, or hybrid energy systems, so research on low-voltage distribution networks with multiple prosumers, as considered in this thesis, remains limited.

In addition, distribution network safety constraints are often simplified in the HEMS and DRL literature. EVs, battery storage systems, and PV systems affect voltage drop, reverse power flow, and transformer loading, while reinforcement learning has been applied to distribution network stability control~\parencite{9960416}. Zhang et al.\ also demonstrated in their secure HEM framework that an unrestricted action space may violate physical constraints, thereby motivating the development of safe MADRL methods~\parencite{ZHANG2024124155}. However, the objective of this thesis is not to propose an entirely new safety projection theory. Instead, this work compares different control methods for multi-prosumer EV--BESS--PV scheduling and analyzes the effects of soft constraints, hard projection, and emergency EV constraint handling on economic performance, EV feasibility, and distribution network safety.

Based on the above research gaps, the research objective of this thesis can be summarized as follows. A multi-prosumer energy management environment is established for a low-voltage distribution network, incorporating household loads, PV generation, stationary BESS, residential EV charging, electricity prices, and power-flow feedback. Rule-based control, Local MPC, ADMM-MPC, and MADRL controllers are compared. Particular emphasis is placed on the soft-constraint and hard-projection approaches for handling the EV departure SoC requirement, together with their impacts on total cost, battery arbitrage, EV charging cost, departure SoC, and distribution network safety. This thesis aims to bridge the gap between single-household HEMS, optimization-based control methods, and MADRL approaches, thereby providing more comprehensive experimental evidence for the coordinated control of EVs and energy storage systems in multi-prosumer scenarios.

